\documentclass[a4paper]{article}
%\usepackage{exercise}
%um nur aufgaben zu zeigen
\usepackage[noanswer]{exercise} 
\usepackage{../images/preamble}
\usepackage{rotating}
\usetikzlibrary{decorations.pathmorphing}
\usetikzlibrary{decorations.markings}
\usetikzlibrary{arrows}
\usetikzlibrary{shapes.geometric}
\newcommand{\midarrow}{\tikz \draw[-triangle 90] (0,0) -- +(.02,0);}
\usepackage{xcolor}
%\usepackage{draftwatermark}
%\SetWatermarkText{\textsc{Draft 2}}
%\SetWatermarkScale{3}
%\SetWatermarkColor{red!30}

\usepackage[printwatermark]{xwatermark}
%\newsavebox\mybox
%\savebox\mybox{\tikz[color=red,opacity=0.3]\node{\textsc{Entwurf}};}
%\newwatermark*[
%allpages,
%angle=45,
%scale=10,
%xpos=-4cm,
%ypos=4cm
%]{\usebox\mybox}
\pagestyle{fancy}
\fancyhead[L]{\includegraphics[width=2cm]{../images/logo_scaled.pdf}}
\fancyhead[R]{\textsc{Lösung Aufgabenserie 5}}


\begin{document}
	\vspace*{-1cm}
	\parbox{4cm}{\vspace{-0.2cm}\includegraphics[width=5cm]{../images/logo_scaled.pdf}}
	\parbox{10.6cm}{\setstretch{2.0} \centering{ \huge \textsf{Aufgabenserie 5 
			}}\\
			Abgabe: 30. November \\ \vspace*{-.5cm} }
		\vspace{0.5cm}

\thispagestyle{empty}
\begin{framed}
	\noindent
	\scriptsize
	Die Aufgaben sollten bis zum \textbf{30. Novermber} bearbeitet werden. Die Lösungen schickt ihr an \href{mailto:physikrolf@gmail.com}{physikrolf@gmail.com}.
	Jede Aufgabe hat eine bestimmte Anzahl an erreichbaren Punkten. Wie viele das sind, müsst ihr raten. Versucht, die Lösungen so genau wie möglich aufzuschreiben. Für besonders schnelle/gute/witzige Lösungen kann es Bonuspunkte geben.\\ Die aktuellen Aufgaben sowie alle alten Aufgabenserien mit Lösungen findet ihr auch auf \url{pankratius.github.io/rolf}. %\\\textit{Zu jeder Aufgabe gibt es jetzt Tipps. Die sollten beim Lösen der Aufgaben helfen.\\ Sollte das so sein macht bitte in euren Lösungen kenntlich, dass bestimmte Schritte von den Tipps und nicht von euch kommen. Darauf gibt es keinen Abzug, es ist nur für uns gut zu wissen.}
\end{framed}

\noindent

\begin{Exercise}[label = voltmeter, difficulty = 3, origin = EstPhO 2003, title = Viele Voltmeter]
	In der abgebildeten Schaltung sind 15 identische Voltmeter und 15 unterschiedliche Ampermeter verbaut, und an eine Batterie angeschlossen. \\
	Das erste Voltmeter zeigt eine Spannung von $U_1 = 9~\mathrm{V}$ an, das erste Ampermeter einen Strom von $I_1 = 2.9~\mathrm{mA}$ und das zweite einen Strom von $I_2 = 2.6~\mathrm{mA}$. Wie groß ist die Summe der Spannungen, die die anderen Voltmeter anzeigen?
\end{Exercise}
\begin{figure}[h]
	\centering
	\tikzset{component/.style={draw,thick,circle,fill=white,minimum size =0.75cm,inner sep=0pt}}
	\begin{tikzpicture}
		\draw (0,1) to[battery1] (0,-1);
		\draw (0,1) to node[component]{$I_1$} (2,1);
		\draw (2,1) to node[component]{$U_1$}(2,-1);
		\draw (2,-1) -- (0,-1);
		\draw (2,1) to node[component]{$I_2$} (4,1);
		\draw (4,1) to node[component]{$V_2$} (4,-1);
		\draw (4,-1) -- (2,-1);
		\draw (4,1) -- (5,1);
		\draw (4,-1) -- (5,-1);
		\node at (5.5,1) {...};
		\node at (5.5,-1) {...};
		
		\draw (6,1) to node[component]{$I_{14}$} (8,1);
		\draw (8,1) to node[component]{$U_{14}$} (8,-1);
		\draw (8,-1) -- (6,-1);
		\draw (8,1) to node[component]{$I_{15}$} (10,1);
		\draw (10,1) to node[component]{$U_{15}$} (10,-1);
		\draw (10,-1) -- (8,-1);
		
		\foreach \n in {2,4,8}{
			\filldraw (\n,1) circle (2pt);
			\filldraw (\n,-1) circle (2pt);
			}
		
		
	\end{tikzpicture}
\end{figure}
\begin{minipage}[b]{0.75\textwidth}
\begin{Exercise}[difficulty = 3, origin = Orpheus e.V, title = Fluß, label = cori1]
	Wir befinden uns auf einer geographischen Breite von $15^\circ$ auf der Nordhalbkugel, und betrachten einen Fluß, der direkt von Norden nach Süden fließt.\\ Der Fluß hat eine Fließgeschwindigkeit von $2~\mathrm{m\cdot s^{-1}}$, und eine Breite von $b = 1000~\mathrm{m}$. Wie groß ist der durch die Erdrotation hervorgerufene Pegelunterschied zwischen dem westlichen und östlichen Flußufer, wenn sich die Flußoberfläche so einstellt, dass sie senkrecht zur wirkenden Gesamtkraft steht. Die Winkelgeschwindigkeit der Erde beträgt $\omega \approx 7.29~\mathrm{s^{-1}}$.
\end{Exercise}
\end{minipage}
\hfill
\begin{minipage}[b]{0.2\textwidth}
	\tikzset{->-/.style={decoration={
			markings,
			mark=at position .5 with {\arrow{>}}},postaction={decorate}}}
\begin{tikzpicture}
\node[draw,ellipse,minimum width=4cm,minimum height=2cm] (ell) {};
\foreach \ang in {-80,-20,-10,0,10,20,80} {
  \pgfmathtruncatemacro{\rang}{180 - \ang}
  \draw[very thin, dashed,shorten >=1pt,shorten <=1pt] (ell.\ang) -- (ell.\rang);
}
\foreach \ang in {-65,-45,...,65} {
  \pgfmathsetmacro{\xrad}{2*sin(\ang)}
  \draw[thin, dashed] (0,1) arc [x radius=\xrad,y radius=1,start angle=90,end angle=-90];
}
\pgfmathsetmacro{\xrad}{2*sin(10)}
\draw[very thick, blue,->-] (0,1) arc[x radius = \xrad, y radius = 1, start angle = 90, end angle = -90];
  \pgfmathtruncatemacro{\rang}{180 - 15}
  \draw[thick, shorten >=1pt,shorten <=1pt,] (ell.15) -- (ell.\rang);
 \node at (2.25,.65) {$15^\circ~\mathrm{N}$};

\end{tikzpicture}
\end{minipage}

\begin{Exercise}[difficulty = 3, origin = {IV. IPhO, 1970}, label = sphmi, title = Teleskop]
	In einem Teleskop wird ein sphärischer Spiegel mit einem Krümmungsradius von $R = 2~\mathrm{m}$ verwendet. Im Brennpunkt des Spiegels befindet sich der Strahlungsempfänger $E$ in Form einer runden Scheibe.  Diese liegt senkrecht zur optischen Achse des Teleskops.\\
	Welche Größe muss der Empfänger haben, damit er die gesamt Strahlung empfängt, die vom Spiegel mit dem Radius von $r = 0,25~\mathrm{m}$ reflektiert wird?
\end{Exercise}



\end{document}